%% Cover letter for the Information Sciences submission.
\documentclass[11pt]{article}
\usepackage[margin=1.15in]{geometry}
\usepackage[T1]{fontenc}
\usepackage{microtype}
\usepackage{xcolor}
\usepackage[colorlinks=true,linkcolor=blue!55!black,urlcolor=blue!55!black]{hyperref}
\usepackage{parskip}
\pagestyle{empty}
\begin{document}

\noindent
Ihor Kendiukhov\\
Institute of Medical Genetics and Applied Genomics\\
University of T\"ubingen, T\"ubingen, Germany\\
\href{mailto:ihor.kendiukhov@student.uni-tuebingen.de}{ihor.kendiukhov@student.uni-tuebingen.de}

\bigskip
\noindent
To the Editors-in-Chief, \emph{Information Sciences}

\bigskip
\noindent
Dear Professor Senatore and Professor Yan,

I am submitting the manuscript \emph{Inference, Not Just Optimisation: Streaming Inversion-Free
Quasi-Newton Estimation on Dependent Data Streams} for consideration as an original research
article in \emph{Information Sciences}.

\textbf{The problem.} Streaming quasi-Newton methods fit a model in a single pass at the cost of
stochastic gradient descent while retaining the efficiency of Newton's method, and they are now
well developed. Their inference is not. The confidence intervals they report are derived assuming
independent observations, and the streams these methods are built for --- sensor networks, traffic,
air quality, click logs --- are serially dependent. The manuscript shows that the consequence is
not a small correction that vanishes with more data: the reported variance omits the
autocovariances of the score, so a nominal 95\% interval converges to a coverage below 95\% and
stays there. We derive that limiting coverage exactly, as a function of a single
variance-inflation factor, and confirm the prediction to within 0.009 across a sweep of dependence
strengths with no fitted constants. On the hourly air-quality stream we study, the limit is near
47\%.

\textbf{The contribution.} The repair is a change of reading order rather than an extra estimation
stage. If the stream is consumed in contiguous blocks with one preconditioned step per block, then
the block-averaged gradients the algorithm already forms are, up to a factor, the non-overlapping
batch means of the score process --- that is, an estimator of precisely the long-run variance the
interval needs. The statistically correct object and the computational shortcut coincide, and the
valid interval ends up \emph{cheaper} than the invalid one it replaces, because an instantaneous
plug-in covariance costs $O(d^2)$ per observation whereas the blocked long-run covariance costs
$O(d^2)$ per block. We prove almost sure convergence, an $O(\log N / N)$ rate and a central limit
theorem with the correct long-run sandwich for the inversion-free rank-one curvature recursion
under geometric mixing --- the regime in which the martingale argument behind the independent-data
analysis fails --- and we give a covariance estimator with an exponentially small block-length
bias whose block length is set by the mixing time of the \emph{data} rather than of the
optimisation path, so that $b \asymp \log N$ suffices.

\textbf{What we report against ourselves.} Blocking the gradient is not free, and the manuscript
says so. It makes the step's contraction condition non-trivial: on a dependent stream a block of
$b$ observations can span far fewer than $b$ distinct covariate directions, and without a
safeguard our own method diverges on a persistent regime-switching design. We state the condition,
give an $O(d)$ remedy, prove it is inactive after finitely many steps, and report the alternative
remedy we implemented first and rejected on measured grounds. Two designs in which no streaming
point estimate has reached its asymptotic regime at the sample sizes we can afford are reported
against an infeasible oracle interval rather than against the nominal level, and on one of them
the residual gap is ours and is traced to a finite-sample bias we show closing. Section~6 lists
what we do not prove, including the two places where a hypothesis is stated but not verified.

\textbf{Fit and reproducibility.} The work sits in the journal's scope under design of algorithms,
symbolic/numeric and statistical techniques, and data engineering: it concerns the algorithmic and
inferential foundations of one-pass estimation on streams. Everything is reproducible: the
estimator, the experiment scripts, the saved results from which every number, figure and table in
the manuscript is generated, and tests that validate every closed-form population quantity against
long simulation are archived with a persistent identifier at
\href{https://doi.org/10.5281/zenodo.21891313}{\texttt{https://doi.org/10.5281/zenodo.21891313}}
(development version:
\url{https://github.com/Kendiukhov/streaming-quasi-newton-dependent-data}), and cited in the
reference list as a software reference. Each of the 100 references was verified against its
published source. The total compute is a few CPU-hours on one machine.

\textbf{Notes on format.} The manuscript is a theoretical contribution and its main narrative is
approximately 42 double-spaced pages, within the 45-page guideline, with the proofs in appendices
as is usual for this kind of work. The main text uses exactly ten figures plus tables, matching
the journal's guideline; four further figures, each supporting a claim the manuscript also
establishes by other means, are supplied as Supplementary Material. I confirm that this work is
original, has not been published elsewhere and is not under consideration by another journal, that
there are no competing interests and no specific funding to declare, and that I have included the
declaration on the use of generative AI tools required by the journal, which I ask you to read: a
large language model was used substantially, under my direction and review, in implementing the
software, running the experiments and drafting the manuscript, and I take full responsibility for
its contents.

Thank you for considering this submission.

\bigskip
\noindent
Yours sincerely,

\bigskip
\noindent
Ihor Kendiukhov

\end{document}
